\section{Determinanter}
Kun defineret for kvadratiske matricer.

\subsection{Determinant af 2×2 og 3×3 matrix}
\begin{align*}
    \begin{bmatrix}
        a
    \end{bmatrix} &= a \\
    \det
    \begin{bmatrix}
        a & b \\
        c & d
    \end{bmatrix}
    &=
    ad - bc \\
    \det
    \begin{bmatrix}
        a & b & c \\
        d & e & f \\
        g & h & i
    \end{bmatrix}
    &=
    aei + bfg + cdh - ceg - bdi - afh
    \kildetag{Sarrus' regel}
\end{align*}

\subsection{Determinant af større matricer (Laplace udvikling)}
\[
\det
\begin{bmatrix}
    a & b & c \\
    d & e & f \\
    g & h & i
\end{bmatrix}
=
a \cdot \det \begin{bmatrix}
    e & f \\
    h & i
\end{bmatrix} - 
b \cdot \det \begin{bmatrix}
    d & f \\
    g & i
\end{bmatrix} +
c \cdot \det \begin{bmatrix}
    d & e \\
    g & h
\end{bmatrix}
\]


General formel for en valgfri $i$'te række:

\[
    \det(\mathbf{A}) = \sum_{j=1}^{n} (-1)^{i+j} \cdot a_{ij} \cdot \det(\mathbf{A}(i;j))
\]

Generel formel for en valgfri $j$'te kolonne:

\[
    \det(\mathbf{A}) = \sum_{i=1}^{n} (-1)^{i+j} \cdot a_{ij} \cdot \det(\mathbf{A}(i;j))
\]

\kilde{Sætning 8.2.1}

\bemærkBig{
$\mathbf{A}(i;j)$ er matricen dannet ved at fjerne den $i$'te række og $j$'te kolonne fra matrix $\mathbf{A}$.

Enhver række eller kolonne kan vælges - Brugbart hvis en række eller kolonne indeholder mange nuller.
}

Man tager altså hvert element i den valgte række eller kolonne, ganger det med determinanten af minoren dannet ved at fjerne den række og kolonne, som elementet befinder sig i, og ganger med $(-1)^{i+j}$ for at tage højde for fortegnet.
\renewcommand{\arraystretch}{1.3}
\begin{center}
    \textbf{Fortegns skema for en $n \times n$ matrix}
    
    \begin{tabular}{|c|c|c|c|}
        \hline
        $+a_{11}$ & $-a_{12}$ & $+a_{13}$ & $\cdots$ \\
        \hline
        $-a_{21}$ & $+a_{22}$ & $-a_{23}$ & $\cdots$ \\
        \hline
        $+a_{31}$ & $-a_{32}$ & $+a_{33}$ & $\cdots$ \\
        \hline
        $\vdots$  & $\vdots$  & $\vdots$  & $\ddots$ \\
    \hline
    \end{tabular}
\end{center}



\subsection{Egenskaber ved determinanter}
Hvis der er en nul-række, er determinanten $0$. 

\kilde{Lemma 8.1.4}

\

Hvis to rækker er ens, er determinanten $0$.

\kilde{Proposition 8.3.5}

\

Hvis to rækker byttes om, skifter determinanten fortegn.

\kilde{Sætning 8.3.6}

\

Hvis en række eller søjle ganges med en skalar $k$, ganges determinanten med $k$.

\kilde{Korollar 8.3.2}

\

Hvis en række lægges til en anden række, ændres determinanten ikke.

\kilde{Sætning 8.3.7}

\

$\det(A) \ne 0 \iff$ det tilsvarende lineære ligningssystem's eneste løsning er nulvektoren. 

\kilde{Korollar 8.2.6}

\begin{align*}
    \det (\mathbf{A}^{-1} ) &= \frac{1}{\det(\mathbf{A})} \kildetag{Sætning 8.1.1} \\
    \text{For diagonalmatricer: } \det(\mathbf{D}) &= \prod_{i=1}^{n} d_{ii} \kildetag{Proposition 8.1.2} \\
    \text{For trekantsmatricer: } \det(\mathbf{T}) &= \prod_{i=1}^{n} t_{ii} \kildetag{Sætning 8.1.(3/5)} \\
    \det(A^T) &= \det(A) \kildetag{Korollar 8.2.2} \\
    \det (\mathbf{A} \cdot \mathbf{B}) &= \det(\mathbf{A}) \cdot \det(\mathbf{B}) \kildetag{Sætning 8.2.3}
\end{align*}